\documentclass[12pt,a4paper]{article}
\usepackage[T2A]{fontenc}
\usepackage[utf8]{inputenc}
\usepackage[russian]{babel}
\usepackage{float}
\usepackage{amsmath, amssymb}
\usepackage{geometry}
\usepackage{tikz}
\usepackage{tikz-3dplot}
\usepackage{caption}
\usetikzlibrary{arrows.meta, calc, decorations.pathmorphing}
\geometry{margin=2cm}

\newcommand{\vr}{\bar r}

\begin{document}

\setcounter{page}{84}

\subsection*{\textbf{ЛЕКЦИЯ 12. КОМПЛЕКСНЫЕ ЧИСЛА}}

\subsection*{\textbf{12.1. Определение комплексного числа.
Алгебраическая форма комплексного числа. Комплексная плоскость}}

\textbf{Определение.} Комплексными числами называется множество выражений вида
$z=x+iy$, где $x,y\in\mathbb R$, а $i$ --- мнимая единица,
удовлетворяющая условию $i^2=-1$, с введенными в нем операциями
сложения и умножения:
$$
z_1=x_1+iy_1;\qquad z_2=x_2+iy_2;
$$
$$
z_1+z_2=(x_1+x_2)+i(y_1+y_2);\qquad
z_1\cdot z_2=(x_1x_2-y_1y_2)+i(x_1y_2+y_1x_2).
$$
Два комплексных числа $z_1=x_1+iy_1$ и $z_2=x_2+iy_2$ равны тогда и
только тогда, когда $x_1=x_2$ и $y_1=y_2$.

$x$ --- действительная часть комплексного числа $z$.

$y$ --- мнимая часть комплексного числа $z$.

Обозначение. $x=\operatorname{Re}z$; $y=\operatorname{Im}z$.

$z=x+iy$ --- алгебраическая форма комплексного числа.

\textbf{Геометрическое изображение комплексных чисел}

Пусть на плоскости задана декартова система координат $Oxy$.
Сопоставим комплексному числу $z=x+iy$ точку $M(x,y)$.

\begin{figure}[H]
  \centering
  \begin{tikzpicture}[scale=1.05, >=Latex]
    \draw[->] (-2.3,0) -- (4.1,0) node[below] {$x$};
    \draw[->] (0,-2.1) -- (0,3.2) node[left] {$y$};
    \draw (1,0.08)--(1,-0.08) node[below] {$1$};
    \draw (-1,0.08)--(-1,-0.08) node[below] {$-1$};
    \draw (0.08,1)--(-0.08,1) node[left] {$i$};
    \draw (0.08,-1)--(-0.08,-1) node[left] {$-i$};
    \node[below left] at (0,0) {$0$};
    \draw[->, thick] (0,0) -- (2.8,2.0) node[above right] {$z$};
    \draw[dashed] (2.8,0) -- (2.8,2.0);
    \draw[dashed] (0,2.0) -- (2.8,2.0);
    \node[below] at (2.8,0) {$x$};
    \node[left] at (0,2.0) {$iy$};
    \node at (2.0,1.15) {$\bar r$};
    \node[below] at (3.0,-0.55) {действительная ось};
    \node[rotate=90] at (-0.9,2.1) {мнимая ось};
    \node at (2.1,2.8) {комплексная плоскость};
  \end{tikzpicture}
  \caption*{Рисунок 12.1}
\end{figure}

\newpage

\textbf{Замечание.} Т.к. геометрически изображать комплексное число
можно не только точкой, но и вектором, то и сложение, вычитание, а
также умножение комплексного числа на действительное число можно
выполнять аналогично тому, как это осуществляется для геометрических векторов.

\textbf{Пример}

\begin{figure}[H]
  \centering
  \begin{tikzpicture}[scale=1.1, >=Latex]
    \draw[->] (-1.4,0) -- (4.5,0) node[below] {$x$};
    \draw[->] (0,-2.2) -- (0,3.5) node[left] {$y$};
    \node[below left] at (0,0) {$0$};
    \draw (1,0.08)--(1,-0.08) node[below] {$1$};
    \draw (-1,0.08)--(-1,-0.08) node[below] {$-1$};
    \draw (0.08,1)--(-0.08,1) node[left] {$i$};
    \draw (0.08,-1)--(-0.08,-1) node[left] {$-i$};
    \coordinate (O) at (0,0);
    \coordinate (A) at (1.1,2.2);
    \coordinate (B) at (3.3,1.0);
    \coordinate (C) at (4.4,3.2);
    \draw[->, thick] (O) -- (A) node[above left] {$z_1$};
    \draw[->, thick] (O) -- (B) node[below right] {$z_2$};
    \draw[->, thick] (O) -- (C) node[above right] {$z_1+z_2$};
    \draw[->] (A) -- (C);
    \draw[->] (B) -- (C);
    \node at (0.45,1.15) {$\bar r_1$};
    \node at (1.7,0.35) {$\bar r_2$};
    \node[rotate=27] at (2.4,1.8) {$\bar r_1+\bar r_2$};
  \end{tikzpicture}
  \caption*{Рисунок 12.2}
\end{figure}

\subsection*{\textbf{12.2. Тригонометрическая форма комплексного
числа. Операции над комплексными числами}}

\begin{figure}[H]
  \centering
  \begin{minipage}{0.42\textwidth}
    \centering
    \begin{tikzpicture}[scale=1.2, >=Latex]
      \draw[->] (-0.2,0) -- (3.5,0) node[below] {$x$};
      \draw[->] (0,-0.2) -- (0,3.0) node[left] {$y$};
      \node[below left] at (0,0) {$0$};
      \coordinate (Z) at (2.5,2.1);
      \draw[->, thick] (0,0) -- (Z) node[above right] {$z$};
      \draw[dashed] (2.5,0) node[below] {$x$} -- (Z);
      \draw[dashed] (0,2.1) node[left] {$iy$} -- (Z);
      \draw (0.7,0) arc (0:40:0.7);
      \node at (0.9,0.35) {$\varphi$};
      \node at (1.3,1.3) {$r$};
    \end{tikzpicture}
    \caption*{Рисунок 12.3}
  \end{minipage}%
  \begin{minipage}{0.48\textwidth}
    $$
    z=x+iy
    \begin{cases}
      x=r\cos\varphi,\\
      y=r\sin\varphi.
    \end{cases}
    $$
  \end{minipage}
\end{figure}

$$
z=r(\cos\varphi+i\sin\varphi) \text{ --- тригонометрическая форма
комплексного числа;}
$$
$$
r=|z|=\sqrt{x^2+y^2} \text{ --- модуль комплексного числа;}
$$
$$
\varphi=\arg z \text{ --- аргумент комплексного числа } z.
$$
$$
z=x+iy=r\left(\frac{x}{r}+i\frac{y}{r}\right)\Rightarrow
\cos\varphi=\frac{x}{r};\quad \sin\varphi=\frac{y}{r}.
$$

\textbf{Пример.} $z=2\sqrt3-2i$
$$
r=|z|=\sqrt{(2\sqrt3)^2+(-2)^2}=4;
$$
$$
z=2\sqrt3-2i=4\left(\frac{2\sqrt3}{4}-i\frac{2}{4}\right)=4\left(\frac{\sqrt3}{2}-i\frac12\right)=
$$
$$
=\left[\varphi=-\frac{\pi}{6}\right]=4\left(\cos\left(-\frac{\pi}{6}\right)+i\sin\left(-\frac{\pi}{6}\right)\right)
\text{ --- тригонометрическая форма.}
$$

\textbf{Умножение комплексных чисел в тригонометрической форме}

Пусть $z_1=r_1(\cos\varphi_1+i\sin\varphi_1)$,
$z_2=r_2(\cos\varphi_2+i\sin\varphi_2)$.

\begin{equation}
  z_1\cdot z_2=r_1\cdot
  r_2\left(\cos(\varphi_1+\varphi_2)+i\sin(\varphi_1+\varphi_2)\right);
  \tag{12.1}
\end{equation}
$$
|z_1\cdot z_2|=|z_1|\cdot |z_2|;
$$
$$
\arg(z_1\cdot z_2)=\arg z_1+\arg z_2.
$$

\textbf{Комплексно сопряженное число}

\begin{figure}[H]
  \centering
  \begin{minipage}{0.45\textwidth}
    $$
    z=x+iy\to \bar z=x-iy.
    $$
    $$
    \bar z \text{ --- комплексно сопряженное по отношению к } z.
    $$
    Если $z\in\mathbb R$ ($\operatorname{Im}z=0$), то $z=\bar z$.
  \end{minipage}%
  \begin{minipage}{0.45\textwidth}
    \centering
    \begin{tikzpicture}[scale=1.05, >=Latex]
      \draw[->] (-0.3,0) -- (3.4,0) node[below] {$x$};
      \draw[->] (0,-2.6) -- (0,2.8) node[left] {$y$};
      \node[below left] at (0,0) {$0$};
      \coordinate (Z) at (2.5,1.8);
      \coordinate (C) at (2.5,-1.8);
      \draw[->, thick] (0,0) -- (Z) node[above right] {$z$};
      \draw[->, thick] (0,0) -- (C) node[below right] {$\bar z$};
      \draw[dashed] (2.5,-1.8) -- (2.5,1.8);
      \node[below] at (2.5,0) {$x$};
      \node[left] at (0,1.8) {$iy$};
      \node[left] at (0,-1.8) {$-iy$};
    \end{tikzpicture}
    \caption*{Рисунок 12.4}
  \end{minipage}
\end{figure}

\textbf{Свойства операции сопряжения}

1. $\overline{z_1\cdot z_2}=\bar z_1\cdot \bar z_2$;

2. $\overline{\left(\frac{z_1}{z_2}\right)}=\frac{\bar z_1}{\bar z_2}$;

3. $z\cdot \bar z=x^2+y^2=|z|^2$.

\textbf{Деление комплексных чисел}
$$
\frac{z_1}{z_2}=\frac{z_1\cdot \bar z_2}{z_2\cdot \bar
z_2}=\frac{z_1\cdot \bar z_2}{|z_2|^2}.
$$

\textbf{Пример}
$$
\frac{1-i}{1+i}=\frac{(1-i)(1-i)}{(1+i)(1-i)}=\frac{(1-i)^2}{1+1}=\frac{1-1-2i}{2}=-i.
$$

\textbf{Деление комплексных чисел в тригонометрической форме}

Пусть $z_1=r_1(\cos\varphi_1+i\sin\varphi_1)$,
$z_2=r_2(\cos\varphi_2+i\sin\varphi_2)$.
\begin{equation}
  \frac{z_1}{z_2}=\frac{r_1}{r_2}\left(\cos(\varphi_1-\varphi_2)+i\sin(\varphi_1-\varphi_2)\right)
  \tag{12.2}
\end{equation}
$$
\left|\frac{z_1}{z_2}\right|=\frac{|z_1|}{|z_2|};
$$
$$
\arg\left(\frac{z_1}{z_2}\right)=\arg z_1-\arg z_2.
$$

\subsection*{\textbf{12.3. Возведение в степень и извлечение корня}}

\textbf{Возведение в степень комплексного числа}

\textbf{Теорема.} Если $z=r(\cos\varphi+i\sin\varphi)$, то
\begin{equation}
  z^n=r^n(\cos n\varphi+i\sin n\varphi), \tag{12.3}
\end{equation}
для $\forall n\in\mathbb Z$.

(12.3) --- формула Муавра.

\textbf{Доказательство}

1) если $n\in\mathbb N$, то формула (12.3) $\Rightarrow$ из (12.1);

2) если $n=0$, то $z^0=1=r^0(\cos0+i\sin0)$;

3) если $n<0$, то пусть $n=-m$, $m\in\mathbb N$.
$$
z^n=z^{-m}=\frac{1}{z^m}=\frac{1}{r^m(\cos m\varphi+i\sin m\varphi)}=
$$
$$
=\frac{1}{r^m}\cdot
\frac{\cos m\varphi-i\sin m\varphi}{(\cos m\varphi+i\sin
m\varphi)(\cos m\varphi-i\sin m\varphi)}=
$$
$$
=r^{-m}(\cos(-m\varphi)+i\sin(-m\varphi))=
$$
$$
=r^n(\cos n\varphi+i\sin n\varphi).
$$
\textbf{ч.т.д.}

\textbf{Пример}

Пусть $z=1+i\sqrt3$. Найти $z^{30}$ --- ?

\textbf{Решение}
$$
|z|=\sqrt{1+3}=2;
$$
$$
z=1+i\sqrt3=2\left(\frac12+i\frac{\sqrt3}{2}\right)=2\left(\cos\frac{\pi}{3}+i\sin\frac{\pi}{3}\right);
$$
$$
z^{30}=2^{30}\left(\cos\frac{\pi}{3}\cdot30+i\sin\frac{\pi}{3}\cdot30\right)=
$$
$$
=2^{30}(\cos10\pi+i\sin10\pi)=2^{30}(\cos0+i\sin0)=2^{30}.
$$

\textbf{Извлечение корня из комплексного числа}

\textbf{Определение.} Пусть $n\in\mathbb N$. Корнем $n$-ой степени из
комплексного числа $z$ называется комплексное число $\alpha$: $\alpha^n=z$.

\textbf{Теорема.} Для ненулевого комплексного числа
$z=r(\cos\varphi+i\sin\varphi)$ существует ровно $n$ различных корней
$\alpha_0,\alpha_1,\alpha_2,\ldots,\alpha_{n-1}$ $n$-ой степени:
$$
\sqrt[n]{z}=\alpha_k=\sqrt[n]{r}\left(\cos\frac{\varphi+2\pi
k}{n}+i\sin\frac{\varphi+2\pi k}{n}\right),
$$
$$
k=0,1,2,\ldots,n-1.
$$

\textbf{Доказательство}

Пусть $\alpha=\rho(\cos\theta+i\sin\theta)\Rightarrow \alpha^n=z\Rightarrow$
$$
\Rightarrow \rho^n(\cos n\theta+i\sin
n\theta)=r(\cos\varphi+i\sin\varphi)\Leftrightarrow
$$
$$
\Leftrightarrow \rho^n=r \text{ и } n\theta=\varphi+2\pi k,\quad k\in\mathbb Z.
$$
$$
\Rightarrow \alpha_k=\sqrt[n]{r}\left(\cos\frac{\varphi+2\pi
k}{n}+i\sin\frac{\varphi+2\pi k}{n}\right) \text{ --- корни }
n\text{-ой степени из числа } z.
$$
Всего $n$ различных корней $\alpha_k$.

$\alpha_0,\alpha_1,\alpha_2,\ldots,\alpha_{n-1}$ --- лежат на
окружности радиуса $R=\sqrt[n]{r}$ в вершинах правильного $n$-угольника

\textbf{ч.т.д.}

\begin{figure}[H]
\centering
\begin{tikzpicture}[scale=1.2, >=Latex]
\draw[->] (-2.4,0) -- (2.6,0) node[below] {$x$};
\draw[->] (0,-1.8) -- (0,2.2) node[left] {$y$};
\draw (0,0) circle (1.45);
\fill (1.35,0.5) circle (1.5pt) node[right] {$\alpha_0$};
\fill (0.1,1.45) circle (1.5pt) node[above] {$\alpha_1$};
\fill (-1.25,-0.72) circle (1.5pt) node[left] {$\alpha_{n-1}$};
\end{tikzpicture}
\caption*{Рисунок 12.5}
\end{figure}

\subsection*{\textbf{12.4. Показательная форма комплексного числа}}

\begin{equation}
e^{i\varphi}=\cos\varphi+i\sin\varphi, \tag{12.4}
\end{equation}
(12.4) --- формула Эйлера.

Из (12.4)
\begin{equation}
\Rightarrow e^{-i\varphi}=\cos\varphi-i\sin\varphi \tag{12.5}
\end{equation}
Из (12.4) и (12.5) $\Rightarrow
\cos\varphi=\frac{e^{i\varphi}+e^{-i\varphi}}{2}$ и
$\sin\varphi=\frac{e^{i\varphi}-e^{-i\varphi}}{2i}$

$z=re^{i\varphi}$ --- показательная форма комплексного числа

$r=|z|$; $\varphi=\arg z$;

$\bar z=re^{-i\varphi}$.

$z^n=r^ne^{in\varphi}$ --- формула Муавра в показательной форме.

$$
\sqrt[n]{z}=\sqrt[n]{r}e^{i\left(\frac{\varphi}{n}+\frac{2\pi
k}{n}\right)};\quad k=0,1,2,\ldots,n-1.
$$

\textbf{Пример}
$$
\sqrt[3]{-1}
$$

\textbf{Решение}
$$
-1=1(\cos\pi+i\sin\pi)
$$
$$
\sqrt[3]{-1}=\sqrt[3]{1}e^{i\left(\frac{\pi}{3}+\frac{2\pi
k}{3}\right)},\quad k=0,1,2
$$
$$
\alpha_0=e^{i\frac{\pi}{3}}=\cos\frac{\pi}{3}+i\sin\frac{\pi}{3}=\frac12+i\frac{\sqrt3}{2}
$$
$$
\alpha_1=e^{i\pi}=\cos\pi+i\sin\pi=-1
$$
$$
\alpha_2=e^{i\frac{5\pi}{3}}=\cos\frac{5\pi}{3}+i\sin\frac{5\pi}{3}=
$$
$$
=\frac12-i\frac{\sqrt3}{2}
$$

\begin{figure}[H]
\centering
\begin{tikzpicture}[scale=1.25, >=Latex]
\draw[->] (-2.1,0) -- (2.3,0) node[below] {$x$};
\draw[->] (0,-1.8) -- (0,1.9) node[left] {$y$};
\draw (0,0) circle (1.25);
\draw (0,0) -- (1.25,0) node[below right] {$1$};
\fill (0.625,1.082) circle (1.5pt) node[above right] {$\alpha_0$};
\fill (-1.25,0) circle (1.5pt) node[left] {$\alpha_1$};
\fill (0.625,-1.082) circle (1.5pt) node[below right] {$\alpha_2$};
\end{tikzpicture}
\caption*{Рисунок 12.6}
\end{figure}

\textbf{Пример}

Пусть $z_1=3-4i$; $z_2=1+2i$. Найти:

1) $z_1\cdot z_2+\bar z_1\cdot \bar z_2$;

2) $\frac{z_1}{z_2}-\frac{\bar z_2}{\bar z_1}$.

\textbf{Решение}

1) $z_1\cdot z_2+\bar z_1\cdot \bar
z_2=(3-4i)(1+2i)+(3+4i)(1-2i)=3+6i-4i+8+3-6i+4i+8=22$.

2)
$$
\frac{z_1}{z_2}-\frac{\bar z_2}{\bar z_1}=\frac{3-4i}{1+2i}-\frac{1-2i}{3+4i}=
$$
$$
=\frac{(3-4i)(3+4i)-(1-2i)(1+2i)}{(1+2i)(3+4i)}=
$$
$$
=\frac{9+16-1-4}{3+4i+6i-8}=\frac{20}{10i-5}=\frac{4}{2i-1}=-\frac{4}{1-2i}=
$$
$$
=-\frac{4(1+2i)}{(1-2i)(1+2i)}=-\frac{4(1+2i)}{1+4}=-\frac45-\frac85i.
$$

\textbf{Пример}

Пусть $z=2\sqrt3+2i$. Найти:

1) модуль $z$;

2) аргумент $z$;

3) тригонометрическую форму комплексного числа $z$;

4) показательную форму комплексного числа $z$;

5) $z^{48}$;

6) $\sqrt[4]{z}$;

7) $z^{4/5}$.

\textbf{Решение}

1) $|z|=\sqrt{12+4}=4$;

2)
$$z=2\sqrt3+2i=4\left(\frac{2\sqrt3}{4}+\frac24i\right)=4\left(\frac{\sqrt3}{2}+\frac12i\right)=$$
$$
=4\left(\cos\frac{\pi}{6}+i\sin\frac{\pi}{6}\right)\Rightarrow
\varphi=\frac{\pi}{6} \text{ --- аргумент;}
$$

3) $$z=4\left(\cos\frac{\pi}{6}+i\sin\frac{\pi}{6}\right) \text{ ---
тригонометрическая форма;}$$

4) $$z=4e^{i\frac{\pi}{6}} \text{ --- показательная форма;}$$

5)
$$z^{48}=4^{48}e^{i\frac{\pi}{6}\cdot48}=2^{96}e^{i8\pi}=2^{96}e^{i0}=2^{96}(\cos0+i\sin0)=2^{96};$$

6) $$\sqrt[4]{z}=\sqrt[4]{4}e^{i\left(\frac{\pi}{24}+\frac{2\pi
k}{4}\right)}=\sqrt2e^{i\left(\frac{\pi}{24}+\frac{\pi
k}{2}\right)},\quad k=0,1,2,3;$$
$$
\alpha_0=\sqrt2e^{i\frac{\pi}{24}};
$$
$$
\alpha_1=\sqrt2e^{i\frac{13\pi}{24}};
$$
$$
\alpha_2=\sqrt2e^{i\frac{25\pi}{24}};
$$
$$
\alpha_3=\sqrt2e^{i\frac{37\pi}{24}}.
$$

\begin{figure}[H]
\centering
\begin{tikzpicture}[scale=1.2, >=Latex]
\draw[->] (-2.0,0) -- (2.3,0) node[below] {$x$};
\draw[->] (0,-2.0) -- (0,2.2) node[left] {$y$};
\draw (0,0) circle (1.35);
\draw (0,0) -- (1.35,0) node[midway, below] {$\sqrt2$};
\foreach \ang/\name/\pos in
{7.5/\alpha_0/right,97.5/\alpha_1/above,187.5/\alpha_2/left,277.5/\alpha_3/below}
{
\coordinate (P) at (\ang:1.35);
\fill (P) circle (1.5pt) node[\pos] {$\name$};
}
\end{tikzpicture}
\caption*{Рисунок 12.7}
\end{figure}

7) $$z^4=4^4e^{i\frac{\pi}{6}\cdot4}=256e^{i\frac{2\pi}{3}}.$$
$$
z^{4/5}=\sqrt[5]{256}e^{i\left(\frac{2\pi}{15}+\frac{2\pi
k}{5}\right)},\quad k=0,1,2,3,4.
$$
$$
\alpha_0=\sqrt[5]{256}e^{i\frac{2\pi}{15}};
$$
$$
\alpha_1=\sqrt[5]{256}e^{i\frac{8\pi}{15}};
$$
$$
\alpha_2=\sqrt[5]{256}e^{i\frac{14\pi}{15}};
$$
$$
\alpha_3=\sqrt[5]{256}e^{i\frac{4\pi}{3}};
$$
$$
\alpha_4=\sqrt[5]{256}e^{i\frac{26\pi}{15}}.
$$

\begin{figure}[H]
\centering
\begin{tikzpicture}[scale=1.2, >=Latex]
\draw[->] (-2.2,0) -- (2.4,0) node[below] {$x$};
\draw[->] (0,-2.1) -- (0,2.2) node[left] {$y$};
\draw (0,0) circle (1.45);
\draw (0,0) -- (1.45,0) node[midway, below] {$\sqrt[5]{256}$};
\foreach \ang/\name/\pos in
{24/\alpha_0/right,96/\alpha_1/above,168/\alpha_2/left,240/\alpha_3/below
left,312/\alpha_4/below right}
{
\coordinate (P) at (\ang:1.45);
\fill (P) circle (1.5pt) node[\pos] {$\name$};
}
\end{tikzpicture}
\caption*{Рисунок 12.8}
\end{figure}

\textbf{Пример}

Найти $z=\left(\frac{1+\sqrt3i}{1-i}\right)^{32}$.

\textbf{Решение}

$z_1=1+\sqrt3i$;

$$
|z_1|=\sqrt{1+3}=2;
$$
$$
z_1=2\left(\frac12+i\frac{\sqrt3}{2}\right)=2\left(\cos\frac{\pi}{3}+i\sin\frac{\pi}{3}\right)=2e^{i\frac{\pi}{3}};
$$
$$
z_2=1-i;
$$
$$
|z_2|=\sqrt{1+1}=\sqrt2;
$$
$$
z_2=\sqrt2\left(\frac{1}{\sqrt2}-i\frac{1}{\sqrt2}\right)=\sqrt2\left(\cos\left(-\frac{\pi}{4}\right)+i\sin\left(-\frac{\pi}{4}\right)\right)=
$$
$$
=\sqrt2e^{-i\frac{\pi}{4}};
$$
$$
\frac{z_1}{z_2}=\frac{2e^{i\frac{\pi}{3}}}{\sqrt2e^{-i\frac{\pi}{4}}}=\sqrt2e^{i\frac{7\pi}{12}};
$$
$$
z=\left(\frac{1+\sqrt3i}{1-i}\right)^{32}=(\sqrt2)^{32}e^{i\frac{7\pi}{12}\cdot32}=2^{16}e^{i\frac{56\pi}{3}}=2^{16}e^{i\frac{2\pi}{3}}.
$$

\end{document}
